\chapter*{Glossário\markboth{Glossário}{}}  % \markboth{}{} é utilizado para corrigir o cabeçalho.
 \label{glossario:latour}
\addcontentsline{toc}{chapter}{Glossário}
\begin{description}

\noindent Este glossário reúne os conceitos de Bruno Latour, Latour e Woolgar, Annemarie Mol e John Law mobilizados ao longo desta tese. As entradas não pretendem fixar definições canônicas; expõem, antes, o uso analítico que cada conceito recebe nos capítulos em que aparece. As referências remetem às edições consultadas.

\bigskip

%% ────────────────────────────────────────────────────────────────────
\section*{Actante}
\addcontentsline{toc}{section}{Actante}
Termo tomado da semiótica greimasiana e reapropriado por Latour para designar qualquer entidade, humana ou não humana, que age ou tem sua ação representada em uma rede. Um actante é definido pragmaticamente: sua identidade emerge da lista de ações e reações que apresenta quando submetido a provas de força, e não de qualidades intrínsecas previamente fixadas \parencite[p.~127]{Latour2011CienciaEmAcao}. Nesta tese, o conceito permite tratar no mesmo plano analítico pesquisadores, GPUs, algoritmos, contratos institucionais, vozes gravadas em enfermarias e o próprio vírus SARS-CoV-2 (Capítulos 1 e 4). A simetria metodológica exigida pelo conceito é seu aspecto mais exigente: não se pode recorrer a fatores sociais ou naturais preestabelecidos para explicar o desfecho de controvérsias; é preciso observar empiricamente quem ou o que age em cada situação.

\bigskip

%% ────────────────────────────────────────────────────────────────────
\section*{Agência dos textos científicos e funil de interesses}
\addcontentsline{toc}{section}{Agência dos textos científicos e funil de interesses}

Law, em \textit{Laboratories and Texts} \parencite{Law1986Laboratories}, argumenta que o texto científico opera como operador de translação que inscreve seus leitores, conduzindo-os ao longo de um caminho especificado. A força de um artigo não reside apenas em sua lógica interna, mas em sua capacidade de mobilizar elementos heterogêneos, dados, instituições, conceitos e especialistas, para impor uma estrutura ao mundo e exercer influência à distância. É por isso que Law descreve o texto como ``a arma secreta da ciência'': transportável, durável e estruturado, o artigo pode circular além do laboratório e agir sobre leitores que nunca estiveram presentes nas práticas que o produziram \parencite{Law1986Laboratories}. O conceito de \textit{funil de interesses} descreve a estrutura pela qual o artigo começa com um problema amplo para atrair o maior público possível e estreita progressivamente o alcance à medida que direciona o leitor em direção a um conjunto particular de pontos. As referências bibliográficas operam como \textit{interessements} elementares: interpõem forças entre o leitor e sua indiferença ao artigo, emprestando a autoridade de outros atores-rede para conduzi-lo adiante. Callon articula a mecânica do \textit{interessement} textual ao demonstrar que o artigo científico, como o laboratório, reúne forças numa combinação que pode escapar ao controle, e que a solução é dispor as palavras em configuração adequada por meio de uma série de \textit{interessements} encadeados \parencite[pp.~71--72]{Callon1986Mapping}. No Capítulo 4, esse vocabulário é aplicado ao artigo do SPIRA: as tabelas de acurácia, as curvas ROC e a seção de materiais e métodos são analisadas como dispositivos que constroem ``realidade'' por encadeamento de inscrições que tornam progressivamente mais difícil ao leitor contestar o resultado sem dominar as mesmas ferramentas que o produziram \parencite{Law1986Laboratories,Callon1986Mapping}.

\bigskip

%% ────────────────────────────────────────────────────────────────────
\section*{Aliados e alistamento de aliados}
\addcontentsline{toc}{section}{Aliados e alistamento de aliados}

Fazer ciência, para Latour, é o processo de recrutar e estabilizar aliados cada vez mais numerosos e heterogêneos: instrumentos, colegas, financiadores, artigos anteriores, dados, instituições e públicos. A força de uma alegação científica depende menos de sua relação com a Natureza do que da extensão da rede que a sustenta: ``aqueles que não alistaram ninguém são ninguém'' \parencite[p.~239]{Latour2011CienciaEmAcao}. Referências bibliográficas funcionam, nessa lógica, como exército de reforços que o autor mobiliza para tornar o custo da discordância proibitivo \parencite[p.~53]{Latour2011CienciaEmAcao}. No Capítulo 2, o próprio processo de construção do referencial teórico desta tese é narrado como alistamento retrospectivo de aliados em resposta a problemas empíricos. Nos Capítulos 3 e 4, o conceito orienta a análise das alianças que Fábio, Cláudio e Marcelo constroem com IBM, FAPESP, hospitais, pacientes e algoritmos para sustentar a pesquisa em IA.

\bigskip

%% ────────────────────────────────────────────────────────────────────
\section*{Caixa-preta}
\addcontentsline{toc}{section}{Caixa-preta}

Quando um fato científico ou um artefato técnico se estabiliza a ponto de não mais suscitar controvérsias sobre seu funcionamento ou conteúdo, diz-se que ele foi posto em caixa-preta. A expressão, oriunda da cibernética, designa o ponto em que uma afirmação ou dispositivo pode ser utilizado como insumo por outros atores sem que seus pressupostos precisem ser reabertos. Latour formula a Regra 1 de método a partir dessa imagem: estudar a ciência em ação significa chegar antes que fatos e máquinas se tenham transformado em caixas-pretas, ou acompanhar as controvérsias que as reabrem \parencite[p.~420]{Latour2011CienciaEmAcao}. No Capítulo 4, o SPIRA é caracterizado como ciência em construção: o artigo científico estabilizou parcialmente os resultados como imutável móvel, mas a falha de generalização para além dos dados pandêmicos impediu que o sistema alcançasse o estatuto de caixa-preta plenamente fechada.

\bigskip

%% ────────────────────────────────────────────────────────────────────
\section*{Centro de cálculo}
\addcontentsline{toc}{section}{Centro de cálculo}

Ponto de acumulação e combinação de inscrições provenientes de locais distantes. Latour mostra que o poder de alguns atores sobre outros resulta de sua capacidade de reunir, comparar e mobilizar representações de fenômenos que ocorrem em outros pontos da rede. O laboratório científico é um centro de cálculo paradigmático: converte entidades do mundo (amostras, medições, observações) em inscrições combináveis que podem ser comparadas, sobrepostas e enviadas a outros centros \parencite{Latour1990Drawing}. No Capítulo 4, o covideiro funciona como centro de cálculo distribuído e precário: enfermarias, smartphones e hospitais convertem vozes humanas em arquivos de áudio que chegam ao IME-USP para transformação em coeficientes MFCC e, finalmente, em classificações binárias de insuficiência respiratória.

\bigskip

%% ────────────────────────────────────────────────────────────────────
\section*{Ciência em ação e ciência pronta (a imagem de Janus)}
\addcontentsline{toc}{section}{Ciência em ação e ciência pronta (a imagem de Janus)}

Latour propõe a imagem do deus romano Janus, que possui dois rostos voltados a direções opostas, para distinguir dois modos de apresentação da ciência. A ciência pronta é a que aparece nos artigos publicados e nos manuais: resultados estabilizados, fatos consolidados, caixas-pretas fechadas, retórica depurada. A ciência em ação é aquela que se observa antes da estabilização: controvérsias abertas, aliados em disputa, instrumentos que falham, dados que não se encaixam, negociações em andamento \parencite[cap.~1]{Latour2011CienciaEmAcao}. O subtítulo desta tese, \textit{como seguir cientistas e engenheiros universidade afora}, é variação direta do subtítulo de \textit{Ciência em ação} e sinaliza a escolha metodológica de observar o segundo rosto de Janus. O SPIRA é analisado no Capítulo 4 precisamente como ciência em construção: o artigo publicado exibe o rosto da ciência pronta (96,5\% de acurácia); a etnografia do covideiro expõe o rosto da ciência em ação (fonoaudiólogos que abandonaram, ruídos que saturaram o sinal, modelos que não generalizaram).

\bigskip

%% ────────────────────────────────────────────────────────────────────
\section*{Controvérsia}
\addcontentsline{toc}{section}{Controvérsia}

Situação em que a validade de um fato científico ou o funcionamento de um artefato técnico ainda está em disputa entre grupos de atores. Para Latour, as controvérsias são as janelas metodológicas por excelência: é nelas que se torna visível o trabalho de construção que a ciência pronta subsequentemente apaga \parencite[cap.~2]{Latour2011CienciaEmAcao}. Acompanhar controvérsias equivale a acompanhar ciência em ação. Na Regra 3, Latour adverte que a resolução de uma controvérsia é a causa da representação da Natureza, e não sua consequência: não se pode invocar a Natureza para explicar por que uma disputa foi encerrada desse modo. No Capítulo 4, a limitação de generalização do SPIRA mantém aberta uma controvérsia técnica: os modelos treinados com dados de COVID-19 não se aplicam bem a insuficiências respiratórias de outras etiologias, e essa abertura é analisada como evidência da dependência do sistema em relação às condições que o produziram.

\bigskip

%% ────────────────────────────────────────────────────────────────────
\section*{Enactment e multiplicidade ontológica}
\addcontentsline{toc}{section}{Enactment e multiplicidade ontológica}

Mol, em \textit{The Body Multiple} \parencite{mol2002body}, propõe uma ruptura com a epistemologia da representação ao demonstrar que práticas distintas não observam o mesmo objeto de perspectivas diferentes, mas produzem objetos ontologicamente distintos. A aterosclerose da sala de cirurgia, do laboratório de patologia e da consulta ambulatorial não são versões de uma mesma doença; são objetos diferentes, \textit{enacted} (performados) por práticas irredutíveis entre si. Como Mol formula: ``realidade não precede as práticas mundanas nas quais interagimos com ela, mas é moldada dentro dessas práticas'' \parencite[p.~6]{mol2002body}. O \textit{enactment} designa esse processo pelo qual práticas produzem objetos sem que haja um objeto singular preexistente a ser representado. A multiplicidade resultante é coexistência de realidades parcialmente conectadas entre si, e não pluralidade de perspectivas sobre uma realidade única. No Capítulo 4, o conceito é mobilizado para descrever a insuficiência respiratória do SPIRA: a IR do covideiro (escuta clínica incorporada), a IR do oxímetro (saturação como número) e a IR do classificador binário (pausa acústica computacionalmente codificada) são enactuações distintas que coexistem sem convergir em um objeto único. A tentativa do SPIRA de estabilizar uma dessas versões como circulável é, nessa leitura, uma operação de translação entre objetos ontologicamente distintos, cuja falha de generalização documenta os limites dessa operação. No Capítulo 1, o conceito orienta a descrição dos dois capítulos analíticos da tese como múltiplos \textit{enactments} do C4AI, produzindo versões do Centro que as análises constroem, e não perspectivas diferentes sobre o mesmo Centro \parencite{mol2002body,MolLaw1994}.

\bigskip

%% ────────────────────────────────────────────────────────────────────
\section*{Fracionalidade}
\addcontentsline{toc}{section}{Fracionalidade}

Law propõe o conceito de \textit{fracionalidade} (\textit{fractionality}) para descrever objetos que resistem tanto à singularidade quanto à pluralidade simples: são ``mais do que um e menos do que muitos'' \parencite[pp.~11--12]{LawHassard1999ANT}. O termo retoma a metáfora matemática do fractal: assim como uma linha fractal ocupa mais de uma dimensão mas menos de duas, o objeto fracional não pode ser arranjado em um espaço topologicamente homogêneo nem como singular nem como múltiplo dentro de um único espaço. Em \textit{Aircraft Stories} \parencite{Law2001Aircraft}, Law acrescenta que narrar objetos sociotécnicos como \textit{projetos}, teleológicos e centralmente coordenados, já é uma performance que apaga as condições fracionais do objeto: a nomeação estabiliza o que a análise deve manter aberto. No Capítulo 4, o SPIRA é descrito como objeto fracional: ao mesmo tempo sistema de IA, projeto de saúde pública, conjunto de práticas de coleta de dados, argumento científico em disputa por legitimidade e rastro de um tempo pandêmico que não se fecha em nenhuma dessas descrições. A própria nomeação ``a rede que Marcelo construiu'', que dá título ao capítulo, é reconhecida ali como escolha de entrada que produz uma versão do objeto, e não como descrição do objeto em sua totalidade \parencite{LawHassard1999ANT,Law2001Aircraft}.

\bigskip

%% ────────────────────────────────────────────────────────────────────
\section*{Hinterland}
\addcontentsline{toc}{section}{Hinterland}

Law toma o termo da geografia: o \textit{hinterland} é o território interior que abastece um porto sem aparecer nas trocas que esse porto realiza. Em termos metodológicos, o \textit{hinterland} de uma descrição é o conjunto de práticas, pressupostos, infraestruturas e recursos que permanecem no fundo sem serem tematizados, mas que são condição de possibilidade do que aparece \parencite[pp.~34--42]{Law2004After}. A diferença em relação à simples ausência está no caráter constitutivo do \textit{hinterland}: ele é o que precisa continuar operando silenciosamente para que o objeto da análise possa existir, e não o que foi excluído da análise. O \textit{hinterland} distingue-se da \textit{ausência manifesta} (o que está fora mas é reconhecidamente necessário ao presente) e da \textit{otherness} (o que o frame analítico não consegue tornar presente sem se transformar): são três modos distintos de estar fora que o \textit{method assemblage} produz simultaneamente. Nesta tese, o conceito opera com recorrência nos Capítulos 3 e 4. No Capítulo 3, as práticas situadas de Marcelo permanecem no \textit{hinterland} enquanto a rede longa de contratos e genealogias corporativas ocupa o primeiro plano. No Capítulo 4, a operação se inverte: a rede longa recua ao \textit{hinterland} sem desaparecer, pois continua sendo condição de possibilidade do SPIRA, enquanto as práticas situadas de Marcelo chegam ao primeiro plano. O \textit{hinterland} estende-se, como Law argumenta, muito além de bancadas de laboratório: alcança conhecimento tácito, software, habilidades de gestão, sistemas de financiamento e agendas políticas, cuja lista é interminável e cujas fronteiras são porosas em todas as direções \parencite[pp.~40--41]{Law2004After}.

\bigskip

%% ────────────────────────────────────────────────────────────────────
\section*{Imutável móvel}
\addcontentsline{toc}{section}{Imutável móvel}

Conceito elaborado por Latour para descrever o tipo de inscrição que permite a acumulação e a circulação de conhecimento científico a longas distâncias. Um imutável móvel é, ao mesmo tempo, móvel (pode ser transportado sem perder sua configuração) e imutável (mantém suas propriedades relacionais ao longo do deslocamento). Mapas, gráficos, tabelas, diagramas e artigos científicos são exemplos paradigmáticos: permitem que um fenômeno ocorrido em um lugar seja comparado com fenômenos de outros lugares no mesmo centro de cálculo \parencite{Latour1990Drawing,Latour2011CienciaEmAcao}. No Capítulo 4, o artigo científico sobre o SPIRA é analisado como imutável móvel: circulou o covideiro como representação estável da insuficiência respiratória detectável por áudio. A falha de generalização revelou, porém, que o imutável era parcialmente mutável fora da configuração pandêmica que o produziu.

\bigskip

%% ────────────────────────────────────────────────────────────────────
\section*{Inscrição e dispositivo de inscrição}
\addcontentsline{toc}{section}{Inscrição e dispositivo de inscrição}

Latour e Woolgar introduzem, em \textit{A vida de laboratório}, o conceito de inscrição para designar qualquer transformação que converte uma entidade material em um traço que pode ser registrado, lido e combinado com outros traços \parencite[p.~45--51]{Latour1997VidaLaboratorio}. Dispositivos de inscrição são os aparatos que realizam essa conversão: espectrofotômetros, centrifugadoras, eletroencefalogramas, questionários. A força analítica do conceito está em evidenciar que o que circula nos laboratórios não são as coisas em si, mas suas representações bidimensionais, combináveis e acumuláveis. Latour retoma e desenvolve o conceito em textos posteriores \parencite{Latour2017Esperanca, Latour1990Drawing}. No Capítulo 4, os dispositivos de inscrição do covideiro são analisados em detalhe: smartphones que convertem vozes em arquivos de áudio, algoritmos que convertem arquivos em coeficientes MFCC, classificadores que convertem coeficientes em diagnósticos binários. Cada etapa é uma inscrição que ganha em combinabilidade ao custo da localidade e materialidade do fenômeno original.

\bigskip

%% ────────────────────────────────────────────────────────────────────
\section*{Matter of fact e matter of concern}
\addcontentsline{toc}{section}{Matter of fact e matter of concern}

Distinção proposta por Latour em \textit{Políticas da natureza} e desenvolvida em textos subsequentes para marcar dois regimes de existência dos objetos científicos. Um \textit{matter of fact} (fato consumado) é um objeto estabilizado, depurado de suas condições de produção, apresentado como dado natural independente das redes que o constituíram. Um \textit{matter of concern} (questão de preocupação) é um objeto que ainda carrega as marcas de sua construção, inseparável das controvérsias, dos aliados e das condições que o produziram \parencite{latour2004politics}. A distinção não é ontológica, mas histórica e política: um \textit{matter of concern} pode tornar-se \textit{matter of fact} à medida que as controvérsias se fecham. No Capítulo 4, a insuficiência respiratória detectada pelo SPIRA é descrita como \textit{matter of concern}: não alcançou o estatuto de \textit{matter of fact} porque permanece inseparável das condições pandêmicas que a produziram como preocupação e objeto de intervenção.

\bigskip

%% ────────────────────────────────────────────────────────────────────
\section*{Method assemblage}
\addcontentsline{toc}{section}{Method assemblage}

Law define \textit{method assemblage} como ``a produção ou \textit{crafting} de um feixe de relações ramificadas que gera presença, ausência manifesta e \textit{otherness}'' \parencite[p.~84]{Law2004After}. Toda pesquisa articula essas três dimensões: presença é o que é trazido aqui e agora, tornado visível, condensado em representações; ausência manifesta é o que está ausente mas reconhecidamente necessário ao presente; \textit{otherness} é o que é necessário à presença mas não pode ou não é tornado manifesto pelo frame analítico. O conceito desloca a questão metodológica de ``descrevi corretamente?'' para ``quais realidades escolho fortalecer?'': se método não reporta algo que já está lá, mas ajuda a produzir realidades, as escolhas metodológicas têm dimensão política constitutiva \parencite[p.~143]{Law2004After}. A \textit{otherness} distingue-se do \textit{hinterland} (ver entrada correspondente): enquanto o \textit{hinterland} é reconhecível e rastreável além do corte, a \textit{otherness} é o que o próprio frame analítico não consegue sequer nomear sem se transformar. Nesta tese, o \textit{method assemblage} opera os dois cortes estratégicos descritos no Capítulo 1: o corte pela extensão, que torna presente a rede longa da parceria IBM-C4AI no Capítulo 3, e o corte pela densidade, que torna presentes as práticas situadas de Marcelo no Capítulo 4. No Capítulo 4, as experiências dos pacientes que gravaram voz no covideiro, seus corpos, suas dores e suas mortes, são \textit{otherness} do SPIRA enquanto dispositivo de inscrição: são condição de possibilidade do sistema, mas o frame computacional que o capítulo analisa não tem como torná-las presentes sem mudar de registro inteiramente \parencite[pp.~83--85]{Law2004After}.

\bigskip

%% ────────────────────────────────────────────────────────────────────
\section*{Política ontológica}
\addcontentsline{toc}{section}{Política ontológica}

Mol e Law formulam a \textit{política ontológica} (\textit{ontological politics}) para nomear as implicações políticas da multiplicidade ontológica: se realidades são múltiplas e performadas por práticas, a questão deixa de ser qual representação está correta e passa a ser quais realidades queremos fortalecer ou enfraquecer \parencite[pp.~74--75]{mol1999ontological}. A descrição e a intervenção tornam-se, nessa perspectiva, modos de fazer realidades: descrever o C4AI de um modo específico fortalece certas versões do Centro enquanto deixa outras sem existência. Latour já abrira essa discussão em \textit{Políticas da Natureza} ao formular a composição progressiva do mundo comum; Mol e Law radicalizam a intuição ao propor multiplicidade ontológica em vez de um mundo comum em construção \parencite{latour2004politics,Law2004After}. Law acrescenta que métodos fazem realidades e políticas: método não reporta algo que já está lá, mas ajuda a produzir o que estuda \parencite[p.~143]{Law2004After}. Nesta tese, o conceito aparece no Capítulo 1 para explicitar as implicações dos cortes analíticos que estruturam os Capítulos 3 e 4: cada \textit{enactment} do C4AI fortalece certas versões do Centro e enfraquece outras, e o reconhecimento disso é simultaneamente metodológico e político \parencite{mol1999ontological,Law2004After}.

\bigskip

%% ────────────────────────────────────────────────────────────────────
\section*{Praxiografia}
\addcontentsline{toc}{section}{Praxiografia}

Mol desenvolve a praxiografia como alternativa metodológica à epistemologia da representação. Onde a epistemologia da representação pergunta como práticas distintas representam um objeto preexistente, a praxiografia acompanha o que as práticas fazem: descreve os objetos que elas produzem, mantendo aberta a multiplicidade sem buscar síntese ou resolução \parencite{mol2002body}. O método praxiográfico exige que o analista permaneça junto às práticas, rastreando como objetos são estabilizados, transformados e multiplicados em cada configuração. A praxiografia distingue-se da etnografia convencional ao recusar a separação entre o que existe e como é conhecido: não há objeto preexistente sobre o qual múltiplas perspectivas incidem; há práticas que, ao realizarem seus objetos, os multiplicam. No Capítulo 1, a tese situa seus dois capítulos analíticos como ``múltiplos \textit{enactments} do Centro no sentido praxiográfico que Mol desenvolve em \textit{The Body Multiple}'': os Capítulos 3 e 4 produzem versões distintas do C4AI ao acompanhar práticas e cortes analíticos distintos, e não perspectivas diferentes sobre o mesmo Centro. No Capítulo 4, a noção é retomada em conjunto com Law \parencite{MolLaw1994} para descrever o que cada \textit{method assemblage} produz como presença, ausência manifesta e \textit{otherness}.

\bigskip

%% ────────────────────────────────────────────────────────────────────
\section*{Referência circulante}
\addcontentsline{toc}{section}{Referência circulante}

Conceito desenvolvido por Latour em \textit{A esperança de Pandora} a partir da análise de uma expedição botânica na Amazônia. A referência circulante descreve o modo pelo qual o conhecimento científico se move ao longo de cadeias de transformação que convertem entidades do mundo em formas progressivamente mais combináveis e portáteis. Em cada etapa da cadeia, perde-se algo da materialidade e localidade do fenômeno original e ganha-se algo em forma, portabilidade e combinabilidade. A referência não é uma relação entre uma proposição e um estado de coisas; é uma cadeia reversível de transformações em que ``os fenômenos são aquilo que circula ao longo da cadeia reversível de transformação, perdendo a cada etapa algumas propriedades a fim de ganhar outras que as tornem compatíveis com os centros de cálculo já instalados'' \parencite[p.~84]{Latour2017Esperanca}. No Capítulo 4, a cadeia de referência circulante do SPIRA é descrita como: som bruto da voz do paciente, arquivo de áudio, coeficientes MFCC, classificação binária de insuficiência respiratória. A condição de reversibilidade dessa cadeia é o que torna o sistema analiticamente interessante.

\bigskip

%% ────────────────────────────────────────────────────────────────────
\section*{Rede tecnocientífica e rede sociotécnica}
\addcontentsline{toc}{section}{Rede tecnocientífica e rede sociotécnica}

Para a Teoria Ator-Rede, a distinção entre o social e o técnico, ou entre a ciência e a tecnologia, não descreve diferenças de natureza, mas é ela mesma um resultado de trabalho de purificação que opera dentro das redes. A expressão \textit{rede tecnocientífica} reúne, no mesmo termo, elementos que as taxonomias convencionais separam: laboratórios, contratos, instrumentos, artigos, financiamentos, políticas e pesquisadores compõem uma mesma trama de associações \parencite{Latour2011CienciaEmAcao}. Nesta tese, \textit{rede tecnocientífica} é o termo geral empregado para descrever essas associações; \textit{rede sociotécnica} aparece quando a ênfase recai sobre a articulação específica entre humanos e não humanos (ver nota 13, Capítulo 1). O Capítulo 3 analisa a rede que Fábio e Cláudio construíram entre USP, IBM e FAPESP; o Capítulo 4 acompanha a rede que Marcelo construiu entre hospitais, pacientes, smartphones e algoritmos para produzir o SPIRA.

\bigskip

%% ────────────────────────────────────────────────────────────────────
\section*{Seguir os atores}
\addcontentsline{toc}{section}{Seguir os atores}

Formulação que sintetiza a orientação metodológica central da Teoria Ator-Rede: em vez de aplicar categorias analíticas prévias ao campo, o pesquisador acompanha os próprios atores onde quer que suas associações os levem, mesmo quando isso implica atravessar fronteiras disciplinares, temporais ou geográficas não antecipadas \parencite[p.~21]{Latour2011CienciaEmAcao}. A Regra 5 de método decorre dessa orientação: deve-se permanecer ``tão indeciso quanto os vários atores que seguimos'' e ``fazer uma lista (não importa se longa e heterogênea) daqueles que realmente trabalham'' \parencite[p.~421]{Latour2011CienciaEmAcao}. No Capítulo 3, foi esse princípio que conduziu a investigação da Bluetalks de fevereiro de 2023 até a máquina de tabulação de Herman Hollerith em 1890: não por decisão metodológica prévia de realizar análise histórica, mas porque as associações mobilizadas pelos próprios atores colocavam questões que a etnografia do presente não conseguia responder.

\bigskip

%% ────────────────────────────────────────────────────────────────────
\section*{Simetria metodológica}
\addcontentsline{toc}{section}{Simetria metodológica}

Princípio epistemológico, originalmente formulado por David Bloor no contexto do Programa Forte em sociologia do conhecimento e radicalizado por Latour, segundo o qual o analista deve tratar simetricamente verdade e erro, sucesso e fracasso, natureza e sociedade, humano e não humano. Para Latour, a simetria generalizada exige que não se use a Natureza para explicar a vitória de uma teoria nem a Sociedade para explicar a derrota de outra: ambas devem ser tratadas como resultados estabilizados de controvérsias, não como seus pontos de partida \parencite[cap.~4]{Latour2011CienciaEmAcao}. Nesta tese, a simetria é simultaneamente princípio metodológico e dificuldade empírica: como a autora reconhece no Capítulo 1, a formação nas ciências sociais induz a buscar o social nos fenômenos e a desconfiar das explicações técnicas, o que contradiz o que os capítulos subsequentes exigem analiticamente.

\bigskip

%% ────────────────────────────────────────────────────────────────────
\section*{Sublata}
\addcontentsline{toc}{section}{Sublata}

Termo empregado por Latour em \textit{A esperança de Pandora} para designar as entidades produzidas pelas operações de inscrição ao longo de uma cadeia de referência circulante. Os dados que chegam ao algoritmo não são os fenômenos do mundo; são \textit{sublata}: coisas realizadas por operações de seleção, gravação, segmentação e transformação que deixaram para trás parte do que existia no ponto de partida \parencite[p.~42]{Latour2017Esperanca}. O conceito salienta que os dados científicos não são extração direta da realidade, mas produtos de trabalho que envolve escolhas, perdas e transformações. No Capítulo 4, o conceito é mobilizado para descrever os coeficientes MFCC que chegam ao classificador do SPIRA: já não são o gemido do paciente, mas entidades produzidas por conversões espectrais que pressupõem toda uma cadeia de decisões técnicas e infraestruturais.

\bigskip

%% ────────────────────────────────────────────────────────────────────
\section*{Tecnociência}
\addcontentsline{toc}{section}{Tecnociência}

Termo que Latour emprega para recusar a separação entre ciência e tecnologia enquanto domínios autônomos com lógicas próprias. A junção não é vocabular; é exigência empírica: o que torna ciência o trabalho dos pesquisadores é precisamente a capacidade de alistar e manter associados elementos heterogêneos, de traduzir interesses de públicos distintos, de recrutar aliados em escalas variadas \parencite{Latour2011CienciaEmAcao}. Separar o momento estritamente técnico do momento político ou social não descreve a prática; apaga o trabalho de mobilização que a constitui. No Capítulo 1, o conceito é mobilizado para interpretar as apresentações de Fábio em eventos como as Bluetalks não como divulgação em vez de fazer ciência, mas como parte constitutiva da prática tecnocientífica: recrutamento de aliados, fortalecimento de rede, tradução de interesses de públicos heterogêneos.

\bigskip

%% ────────────────────────────────────────────────────────────────────
\section*{Topologias: fluido, rede e região}
\addcontentsline{toc}{section}{Topologias: fluido, rede e região}

Law, retomando e desenvolvendo argumento elaborado com Mol \parencite{MolLaw1994}, propõe pensar o mundo sociotécnico em termos de \textit{complexidade topológica}: o espaço sociotécnico não é topologicamente uniforme, mas habitado por diferentes formas de espacialidade que coexistem sem se conformar umas às outras. Redes, regiões e fluidos são exemplos de topologias distintas que podem estar simultaneamente em jogo num mesmo fenômeno \parencite{Law1999Actor}. Uma \textit{rede} existe pela estabilidade de suas conexões e muda de forma quando uma conexão se rompe. Uma \textit{região} existe pela continuidade espacial e mantém suas propriedades pela proximidade. Um \textit{fluido} mantém identidade apesar de mudar de forma, adaptando-se ao contexto sem por isso fragmentar-se: persiste pela semelhança de substância e comportamento, e não pela fixidez de conexões. No Capítulo 4, o conceito é operacionalizado para descrever a tensão analítica do SPIRA: o escutar clínico incorporado do médico pertence a uma topologia fluida (situado, incorporado, dependente de quem ouve, resistente à transportação), ao passo que o classificador binário treinado com coeficientes MFCC pertence a uma topologia de rede (transportável, replicável, independente do ouvido situado). O SPIRA é descrito como esforço de \textit{translação inter-topológica}: a tentativa de converter o fluido em rede, e a falha de generalização dos modelos é evidência analítica de que essa conversão permaneceu incompleta. O paralelo com o estudo de Mol sobre a aterosclerose é explicitado no capítulo: enquanto Mol demonstra a coexistência irredutível de múltiplos objetos sem convergência possível, o SPIRA tenta produzir essa convergência tecnicamente, e a análise acompanha o que se perde, o que resiste e o que coexiste nessa tentativa \parencite{MolLaw1994,Law1999Actor,Law2004After}.

\bigskip

%% ────────────────────────────────────────────────────────────────────
\section*{Translação}
\addcontentsline{toc}{section}{Translação}

Processo pelo qual um ator converte, desvia ou redefine os interesses de outros atores de modo a alinhá-los aos seus próprios objetivos. A translação não é um ato isolado, mas uma série de operações que constroem equivalências entre entidades distintas: um problema de saúde pública é transladado em problema computacional; um fenômeno físico é transladado em coeficiente matemático; um interesse corporativo é transladado em agenda de pesquisa universitária \parencite{Latour2011CienciaEmAcao}. Callon articula o conceito com as noções de problematização, interessamento, alistamento e mobilização de porta-vozes. Nesta tese, o conceito de translação aparece no Capítulo 4 na análise do primeiro movimento condicionante do SPIRA: a tradução inter-topológica que converte o escutar clínico incorporado (topologia fluida) em inscrições transportáveis e independentes do ouvido situado (topologia de rede). A falha de generalização dos modelos é, nessa leitura, evidência de que a translação inter-topológica permaneceu incompleta.

\bigskip

%% ────────────────────────────────────────────────────────────────────
\section*{Vida de laboratório e construção de fatos científicos}
\addcontentsline{toc}{section}{Vida de laboratório e construção de fatos científicos}

Em \textit{A vida de laboratório}, Latour e Woolgar realizam etnografia do Salk Institute of Biological Studies, descrevendo como fatos científicos são produzidos a partir de atividades sociais e materiais cotidianas \parencite{Latour1997VidaLaboratorio}. A tese central é que os fatos não existem como tais antes das práticas laboratoriais que os produzem: são construções que passam por estágios de estabilização progressiva, dos enunciados mais contingentes e contextualizados aos mais estabilizados e descontextualizados. Essa etnografia fundadora demonstra que o laboratório é o lugar onde inscrições são produzidas, combinadas e convertidas em afirmações sobre o mundo. Nesta tese, a orientação metodológica de \textit{A vida de laboratório} é operacionalizada em dois contextos distintos: no Capítulo 3, o C4AI é analisado como laboratório institucional onde parcerias constroem condições de possibilidade para a pesquisa em IA; no Capítulo 4, o covideiro é descrito como laboratório distribuído e precário onde vozes humanas são convertidas em dados computacionais.

\bigskip

%% ────────────────────────────────────────────────────────────────────
%% REFERÊNCIAS DESTA SEÇÃO — apenas as utilizadas no glossário
%% (serão integradas à lista geral de referências do trabalho)
%% ────────────────────────────────────────────────────────────────────

% \parencite{Latour2011CienciaEmAcao}   → Ciência em ação (2011)
% \parencite{Latour2001Esperanca}       → A esperança de Pandora (2001)
% \parencite{Latour1990Drawing}         → Drawing things together (1990)
% \parencite{Latour2004Politics}        → Politics of Nature (2004)
% \parencite{LatourWoolgar1997VidaLaboratorio} → A vida de laboratório (1997)
% \parencite{mol2002body}               → The Body Multiple (2002)
% \parencite{mol1999ontological}        → Ontological politics (1999)
% \parencite{MolLaw1994}                → Regions, Networks and Fluids (1994)
% \parencite{Law1999Actor}              → After ANT (1999)
% \parencite{Law2004After}              → After Method (2004)
% \parencite{Law2001Aircraft}           → Aircraft Stories (2001)
% \parencite{LawHassard1999ANT}         → Actor Network Theory and After (1999)
% \parencite{Law1986Laboratories}       → Laboratories and Texts (1986)
% \parencite{Callon1986Mapping}         → Mapping the Dynamics... (1986)
\end{description}